
\documentclass{article}
\usepackage{colortbl}
\usepackage{makecell}
\usepackage{multirow}
\usepackage{supertabular}

\begin{document}

\newcounter{utterance}

\centering \large Interaction Transcript for game `cladder', experiment `full\_v1.5\_default', episode 403 with score{-}Qwen3.5{-}9B{-}sp.
\vspace{24pt}

{ \footnotesize  \setcounter{utterance}{1}
\setlength{\tabcolsep}{0pt}
\begin{supertabular}{c@{$\;$}|p{.15\linewidth}@{}p{.15\linewidth}p{.15\linewidth}p{.15\linewidth}p{.15\linewidth}p{.15\linewidth}}
    \# & \multicolumn{2}{c}{Player} && \multicolumn{2}{c}{Game Master} \\
    \hline

    \theutterance \stepcounter{utterance}  
    & & & \multicolumn{4}{p{0.6\linewidth}}{
        \cellcolor[rgb]{0.9,0.9,0.9}{
            \makecell[{{p{\linewidth}}}]{
                \texttt{\tiny{[P1$\langle$GM]}}
                \texttt{You are an expert in causal inference. The following question is not a typical commonsense query, but rather a meticulously designed question created by a professor specializing in causal inference, intended to assess the students' mastery of the course content.} \\
\texttt{Imagine a self{-}contained, hypothetical world with only the following conditions, and without any unmentioned factors or causal relationships: Vaccination status has a direct effect on vaccination reaction and getting smallpox. Getting smallpox has a direct effect on smallpox survival. Vaccination reaction has a direct effect on smallpox survival. For unvaccinated individuals, the probability of smallpox survival is 39\%. For vaccinated individuals, the probability of smallpox survival is 46\%. For vaccinated individuals, would it be more likely to see smallpox survival if the person had been vaccinated?} \\
\texttt{Start your answer with 'Yes' or 'No', followed by additional reasoning or evidence to support your explanation.} \\
            }
        }
    }
     \\ \\

    \theutterance \stepcounter{utterance}  
    & \multicolumn{4}{p{0.6\linewidth}}{
        \cellcolor[rgb]{0.561,0.737,0.561}{
            \makecell[{{p{\linewidth}}}]{
                \texttt{\tiny{[P1$\rangle$GM]}}
                \texttt{No} \\
\\ 
\texttt{**Reasoning:**} \\
\\ 
\texttt{To answer this question, we must distinguish between the **marginal probability** of survival (observed in the population) and the **counterfactual probability** of survival (what would have happened under a different intervention).} \\
\\ 
\texttt{1.  **Identify the Observed Data:**} \\
\texttt{    *   \$P(\textbackslash{}text\{Survival\} | \textbackslash{}text\{Vaccinated\}) = 46\textbackslash{}\%\$. This is the probability of survival for individuals who were actually vaccinated.} \\
\\ 
\texttt{2.  **Analyze the Causal Structure:**} \\
\texttt{    *   The prompt states: "Vaccination status has a direct effect on vaccination reaction and getting smallpox."} \\
\texttt{    *   Crucially, it implies that vaccination status influences the likelihood of *getting* smallpox. In standard causal models for vaccination, the primary mechanism of protection is preventing infection (getting smallpox).} \\
\texttt{    *   The prompt also states: "Vaccination reaction has a direct effect on smallpox survival." This suggests that among those who *do* get smallpox, the reaction to the vaccine (or the vaccine's presence) affects survival.} \\
\\ 
\texttt{3.  **Evaluate the Counterfactual (\$P(\textbackslash{}text\{Survival\} | \textbackslash{}text\{Vaccinated\}, \textbackslash{}text\{Do\}(\textbackslash{}text\{Unvaccinated\}))\$):**} \\
\texttt{    *   The question asks: "For vaccinated individuals, would it be more likely to see} \\
            }
        }
    }
    & & \\ \\

    \theutterance \stepcounter{utterance}  
    & & & \multicolumn{2}{p{0.3\linewidth}}{
        \cellcolor[rgb]{0.9,0.9,0.9}{
            \makecell[{{p{\linewidth}}}]{
                \texttt{\tiny{[GM$|$GM]}}
                \texttt{No} \\
            }
        }
    }
    & & \\ \\

    \theutterance \stepcounter{utterance}  
    & & & \multicolumn{2}{p{0.3\linewidth}}{
        \cellcolor[rgb]{0.9,0.9,0.9}{
            \makecell[{{p{\linewidth}}}]{
                \texttt{\tiny{[GM$|$GM]}}
                \texttt{no} \\
            }
        }
    }
    & & \\ \\

    \theutterance \stepcounter{utterance}  
    & & & \multicolumn{2}{p{0.3\linewidth}}{
        \cellcolor[rgb]{0.9,0.9,0.9}{
            \makecell[{{p{\linewidth}}}]{
                \texttt{\tiny{[GM$|$GM]}}
                \texttt{game\_result = WIN} \\
            }
        }
    }
    & & \\ \\

\end{supertabular}
}

\end{document}
